% SET ARTICLE STYLE
\documentclass[paper=a4,fontsize=11pt,twocolumn]{article}


% PACKAGES
\usepackage[english]{babel} % English language/hyphenation
%\usepackage{fontspec} % To use Avenir
\usepackage{geometry} % For paper size and layout
\usepackage{color}
%\usepackage{xcolor}
\usepackage[svgnames]{xcolor} % Enabling colors by their 'svgnames'
\usepackage{fancyhdr} % For header and footer with odd and even page params
\usepackage{lipsum}  % For generating sample text
\usepackage{ragged2e} % For restating fully justified environment
\usepackage{textcomp} % ??? Can't remember...
\usepackage{titlesec} % To edit title formats
\usepackage{lastpage} % Ffor Page X of Y
\usepackage{booktabs} % For nicer tables
\usepackage{multicol} % More things for tables (and arranging text)
\usepackage{graphicx} % For easier figure wrangling
\usepackage{lineno} % Line numbers for review purposes
\usepackage{microtype} % For uneven columns
\usepackage{vwcol} % For uneven columns
\usepackage[T1]{fontenc} % For fonts
\usepackage{lmodern} % {tgheros lmodern paratype} For fonts
\renewcommand{\familydefault}{\sfdefault}

% DEFINITIONS
\linenumbers
 \geometry{
	 a4paper,
	 top=20mm,
	 right=15mm,
	 left=15mm,
 	 bottom=20mm,
 	 }
\definecolor{brand-primary}{HTML}{1A2530}
\definecolor{brand-secondary}{HTML}{2C5A8F}
\definecolor{brand-logo-dark}{HTML}{58778C}
\definecolor{brand-logo-light}{HTML}{BCD3DB}
\definecolor{brand-midlight}{HTML}{C2C5C3}

\titleformat*{\section}{\Large\normalfont}
\titleformat*{\subsection}{\large\normalfont}
\titleformat*{\subsubsection}{\large\normalfont}
\titleformat*{\paragraph}{\normalfont}


% HEADERS AND FOOTERS 
\pagestyle{fancy}
\fancyhf{}
\fancyhead[RE,LO]{ARTICLES}

\fancyhead[LE,RO]{\footnotesize\textcolor{brand-secondary}{Badgers and being the best thing ever}}
\fancyfoot[LE,LO]{\footnotesize{\textcopyright \textcolor{brand-secondary}{Journal of the National Student Association of Medical Research}\\
\hspace{4mm}\footnotesize2017 Vol 1 Issue 1 \emph{doi:xx.xxxx/xxx}}}
\fancyfoot[RE,RO]{\footnotesize Page \thepage\ of \pageref{LastPage}}
\renewcommand{\headrulewidth}{2pt}
\renewcommand{\footrulewidth}{1pt}


% MISC
\graphicspath{{figures/}}
\hbadness=100000 % Turn off those dreadful overfull and underful hboxes
\hfuzz=1000pt
\vbadness=100000
\vfuzz=1000pt


% MACROS
\newcommand{\horizontalstrip}[1]{\par\noindent\colorbox{shadecolor}
{\parbox{\dimexpr\textwidth-2\fboxsep\relax}{#1}}}



% START OF DOCUMENT
\begin{document}
\pagestyle{fancy}

\twocolumn[
	\begin{@twocolumnfalse}

		% Title 
		\huge\textcolor{brand-secondary}{Title of this paper: honestly it probably has something to do with badgers}\\	
		
		% Author(s)
		\Large{First X Author$^{*}$, Second X Author$^{**}$ and Third X Author$^{***}$}
		
		% Affiliation(s)
		\setlength{\columnseprule}{1pt}
		\def\columnseprulecolor{\color{brand-secondary}}
		\begin{multicols}{2}
			\small{$^{*\alpha}$affiliation\smallskip\\
				$^{**}$affiliation\smallskip\\
				$^{***}$affiliation}\bigskip\smallskip\\
			\vfill
			\small{$^{\alpha}$Corresponding author: email@domain.ac.uk}
			\columnbreak

		% Dates	
			\hangindent=0.8cm 
			\small{
				\hspace{8mm}Received: 16$^{th}$ January 2017\smallskip\\
				Revised: 31$^{st}$ January 2017, 10$^{th}$ February 2017\smallskip\\
				Accepted: 11$^{th}$ February 2017\smallskip\\
				Available online: 12$^{th}$ February 2017\bigskip\\
				\vfill
				\hspace{8mm}Keywords: \colorbox{brand-midlight}{\strut sample paper} \colorbox{brand-midlight}{\strut badgers} \colorbox{brand-midlight}{\strut awesome}
				}
		\end{multicols}
		\bigskip		

		% Abstract
		\begin{minipage}{0.95\textwidth}
		\begin{centering}
		\setlength{\fboxsep}{10pt}
		\fcolorbox{brand-primary}{brand-logo-light}{\parbox{\textwidth}{
			\Large{Abstract}\bigskip\\
			\normalsize
			\textbf{Introduction}
				words words words badger words words words badger words words words badger words words words badger words words words badger\\
			\textbf{Method}
				words words words badger words words words badger words words words badger words words words badger words words words badger words words words badger words words words badger\\
			\textbf{Results}
				words words words badger words words words badger words words words badger words words words badger words words words badger words words words badger words words words badger words words words badger words words words badger\\
			\textbf{Conclusions}
				words words words badger words words words badger words words words badger words words words badger words words words badger\\
			}
		} % End fbox
		\end{centering}
		\end{minipage} % END ABSTRACT

		
	\end{@twocolumnfalse}
\vspace{8mm}
] % End @twocolumnfalse


% Article text
\section{Introduction}
\lipsum[1-7]

\section{Method}
\subsection{Study design}
Figure fig:one-column shows a one-column figure.

\begin{figure}
\centering
\includegraphics[width=\columnwidth]{logo-long-jnsamr.pdf}
\caption{This is a single-column figure. This figure also demonstrates why the figures need to be vector graphics if at all possible, or, if not, the resolution needs to be reallllly high.}
\label{fig:one-column}
\end{figure}

\lipsum[8-9]

\subsection{Statistical analysis}
\lipsum[10-11]

\section{Results}
\subsection{Effects of X}
\subsubsection{Raw result}

Table \ref{tab:simple} shows an example of a simple table.

\begin{table}
\begin{centering}
\caption{A single-column table}\label{tab:one-column}
	\begin{tabular}[\columnwidth]{ccl}\toprule
		Condition &  Filter cutoff (Hz) &  Further processing \\
		\midrule
		1  & 21591--22050 & 16-bit quantization \\
		2  & 21591--22050 & 16-bit quantization \& \\
		 & & \hspace{8.1pt}RPDF dither \\
		3  & 21591--22050 & None \\
		4  & 23500--24000 & 16-bit quantization \\
		5  & 23500--24000 & 16-bit quantization \& \\
		 & & \hspace{8.1pt}RPDF dither \\
		6  & 23500--24000 & None \\
		\bottomrule
	\end{tabular}\par
	\medskip
\textit{Notes for this table if needed}
\end{centering}
\end{table}

Table \ref{tab:larger} shows an example of a table that requires two columns.

\begin{table*}
\begin{centering}
\caption{A two-column table}\label{tab:two-column}%
	\begin{tabular}[\columnwidth]{ccccccccc}\toprule
	Condition & L1 & L2 & L3 & L4 & L5 & L6 & L7 & L8 \\
	\midrule
	1 & 0.32 & --0.38 & 0.23 & 0.05 & --0.15 & --0.35 & 0.05 & --0.44 \\
	2 & 0.27 & --0.15 & 1.15 & --0.21 & 1.67 & --0.33 & --0.14 & --0.18 \\
	3 & 0.60 & --0.25 & 0.31 & --0.09 & --0.24 & --0.39 & --0.14 & --0.11 \\
	4 & 0.02 & --0.31 & 0.32 & --0.47 & --0.39 & --0.09 & 0.13 & --0.09 \\
	5 & 0.18 & --0.38 & 0.44 & --0.40 & 0.13 & 0.14 & --0.44 & --0.06 \\
	6 & --0.18 & --0.67 & --0.16 & 0.19 & --0.24 & --0.27 & --0.03 & 0.14 \\
	\cmidrule{2-9}
	Mean & 0.20 & --0.36 & 0.38 & --0.15 & 0.13 & --0.21 & --0.09 & --0.12 \\
	\bottomrule
	\end{tabular}\par
	\medskip
\textit{Notes for this table if needed}
\end{centering}
\end{table*}

\lipsum[12]
\subsubsection{After correction for multiple comparisons}
\lipsum[13]

\subsection{Effects of Y}
\subsubsection{Raw result}
\lipsum[14]
\subsubsection{After correction for multiple comparisons}
\lipsum[15]

\subsection{Summary of results}
\lipsum[16-18]

\section{Discussion}
\lipsum[19-23]

\section{Conclusions}
\lipsum[24-32]

\end{document}